\part{I–II klasė}

\section{Skaičiai iki 100 ir veiksmai}\label{sec:skaiciai-iki-100-ir-veiksmai}

\subsection{Apibrėžtys}

\begin{apibreztis}
Skaičius iki 100 parodo, kiek daiktų yra rinkinyje, arba kurią vietą daiktas užima eilėje. Pavyzdžiui, skaičius 37 reiškia 3 dešimtis ir 7 vienetus.
\end{apibreztis}

\begin{apibreztis}
Dešimtis yra 10 vienetų. Skaičiuje 54 yra 5 dešimtys ir 4 vienetai.
\end{apibreztis}

\begin{apibreztis}
Veiksmas yra skaičių keitimas pagal taisyklę. Šioje temoje svarbiausi veiksmai yra sudėtis ir atimtis.
\end{apibreztis}

\subsection{Teoremos ir savybės}

\begin{savybe}
Skaičiai didėja einant skaičių juosta į dešinę ir mažėja einant į kairę.
\end{savybe}

\begin{savybe}
Skaičių galima išskaidyti į dešimtis ir vienetus: $68=60+8$.
\end{savybe}

\begin{savybe}
Pridėjus 1 gaunamas kitas skaičius, o atėmus 1 gaunamas ankstesnis skaičius.
\end{savybe}

\begin{pastaba}
Pirmoje ir antroje klasėje skaičius patogu matyti kaip daiktus, pirštus, taškus, lazdeles arba vietą skaičių juostoje.
\end{pastaba}

\subsection{Taisyklės ir algoritmai}

\begin{enumerate}
\item Skaičiuojant pirmyn, prie skaičiaus vis pridedame po 1.
\item Skaičiuojant atgal, iš skaičiaus vis atimame po 1.
\item Lyginant du skaičius, pirmiausia žiūrime į dešimtis, tada į vienetus.
\item Norint sudėti arba atimti iki 100, patogu skaičių išskaidyti į dešimtis ir vienetus.
\end{enumerate}

\subsection{Iliustruojantys pavyzdžiai}

\textbf{1 pavyzdys.} Kiek dešimčių ir vienetų yra skaičiuje 46?

Skaičių 46 sudaro 4 dešimtys ir 6 vienetai, nes $46=40+6$.

\textbf{Atsakymas: 4 dešimtys ir 6 vienetai.}

\textbf{2 pavyzdys.} Kuris skaičius didesnis: 58 ar 63?

Skaičiuje 58 yra 5 dešimtys, o skaičiuje 63 yra 6 dešimtys. Kadangi 6 dešimtys yra daugiau negu 5 dešimtys, tai $63>58$.

\textbf{Atsakymas: 63 yra didesnis.}

\textbf{3 pavyzdys.} Apskaičiuok: $34+5$.

Skaičiuojame nuo 34 penkis žingsnius pirmyn: 35, 36, 37, 38, 39. Todėl $34+5=39$.

\textbf{Atsakymas: 39.}

\subsection{Ryšys su kitomis temomis}

Skaičių iki 100 supratimas padeda mokytis sudėties ir atimties temoje \ref{sec:sudetis-ir-atimtis}. Vėliau jis plečiamas iki skaičių iki 1000 temoje \ref{sec:skaiciai-iki-1000-ir-vietine-verte} ir iki didesnių natūraliųjų skaičių temoje \ref{sec:naturalieji-skaiciai-iki-milijono}.

\section{Skaičiai iki 1000 ir vietinė vertė}\label{sec:skaiciai-iki-1000-ir-vietine-verte}

\subsection{Apibrėžtys}

\begin{apibreztis}
Skaičiai iki 1000 yra skaičiai nuo 0 iki 1000. Juos galima sudaryti iš šimtų, dešimčių ir vienetų.
\end{apibreztis}

\begin{apibreztis}
Šimtas yra 10 dešimčių arba 100 vienetų.
\end{apibreztis}

\begin{apibreztis}
Vietinė vertė rodo, ką reiškia skaitmuo pagal savo vietą skaičiuje. Skaičiuje 426 skaitmuo 4 reiškia 4 šimtus, 2 reiškia 2 dešimtis, o 6 reiškia 6 vienetus.
\end{apibreztis}

\subsection{Teoremos ir savybės}

\begin{savybe}
Tas pats skaitmuo skirtingose vietose turi skirtingą vertę. Skaičiuje 707 pirmasis 7 reiškia 700, o paskutinis 7 reiškia 7.
\end{savybe}

\begin{savybe}
Trijų skaitmenų skaičių galima išskaidyti į šimtus, dešimtis ir vienetus: $538=500+30+8$.
\end{savybe}

\begin{savybe}
Jeigu lyginame skaičius iki 1000, pirmiausia lyginame šimtus, tada dešimtis, tada vienetus.
\end{savybe}

\subsection{Taisyklės ir algoritmai}

\begin{enumerate}
\item Perskaityk skaičių iš kairės į dešinę: šimtai, dešimtys, vienetai.
\item Išskaidyk skaičių: šimtai $+$ dešimtys $+$ vienetai.
\item Lygindamas du skaičius, pradėk nuo didžiausios vietos, tai yra nuo šimtų.
\item Apvalindamas iki dešimčių, žiūrėk į vienetus; apvalindamas iki šimtų, žiūrėk į dešimtis.
\end{enumerate}

\subsection{Iliustruojantys pavyzdžiai}

\textbf{1 pavyzdys.} Išskaidyk skaičių 372.

Skaičiuje 372 yra 3 šimtai, 7 dešimtys ir 2 vienetai. Todėl $372=300+70+2$.

\textbf{Atsakymas: $372=300+70+2$.}

\textbf{2 pavyzdys.} Kuris skaičius didesnis: 418 ar 481?

Abu skaičiai turi po 4 šimtus. Lyginame dešimtis: 418 turi 1 dešimtį, o 481 turi 8 dešimtis. Kadangi 8 dešimtys daugiau, tai $481>418$.

\textbf{Atsakymas: 481 yra didesnis.}

\textbf{3 pavyzdys.} Koks skaičius turi 6 šimtus, 0 dešimčių ir 9 vienetus?

6 šimtai yra 600, 0 dešimčių yra 0, 9 vienetai yra 9. Sudedame: $600+0+9=609$.

\textbf{Atsakymas: 609.}

\subsection{Ryšys su kitomis temomis}

Vietinė vertė remiasi skaičiais iki 100 iš temos \ref{sec:skaiciai-iki-100-ir-veiksmai}. Ji reikalinga sudėčiai ir atimčiai \ref{sec:sudetis-ir-atimtis}, vėliau padeda suprasti natūraliuosius skaičius iki milijono temoje \ref{sec:naturalieji-skaiciai-iki-milijono}.

\section{Sudėtis ir atimtis}\label{sec:sudetis-ir-atimtis}

\subsection{Apibrėžtys}

\begin{apibreztis}
Sudėtis yra veiksmas, kai kelios dalys sujungiamos į vieną visumą. Veiksmo ženklas yra $+$.
\end{apibreztis}

\begin{apibreztis}
Atimtis yra veiksmas, kai iš visumos paimama dalis arba lyginamas skirtumas. Veiksmo ženklas yra $-$.
\end{apibreztis}

\begin{apibreztis}
Suma yra sudėties rezultatas, o skirtumas yra atimties rezultatas.
\end{apibreztis}

\subsection{Teoremos ir savybės}

\begin{savybe}
Sudėties dėmenis galima sukeisti vietomis: $4+7=7+4$.
\end{savybe}

\begin{savybe}
Atimtis tikrina sudėtį: jeigu $12-5=7$, tai $7+5=12$.
\end{savybe}

\begin{savybe}
Pridėjus 0 skaičius nepasikeičia, o atėmus 0 skaičius taip pat nepasikeičia.
\end{savybe}

\begin{pastaba}
Atimties dėmenų sukeisti vietomis paprastai negalima: $9-4$ ir $4-9$ pirmose klasėse reiškia skirtingas situacijas.
\end{pastaba}

\subsection{Taisyklės ir algoritmai}

\begin{enumerate}
\item Sudėdamas be perėjimo per dešimtį, sudėk vienetus su vienetais ir dešimtis su dešimtimis.
\item Sudėdamas su perėjimu per dešimtį, pirmiausia papildyk iki artimiausios dešimties.
\item Atimdamas, gali skaičiuoti atgal arba ieškoti, kiek trūksta iki didesnio skaičiaus.
\item Atsakymą patikrink priešingu veiksmu.
\end{enumerate}

\subsection{Iliustruojantys pavyzdžiai}

\textbf{1 pavyzdys.} Apskaičiuok: $28+14$.

Išskaidome: $28+14=28+10+4$. Pirmiausia $28+10=38$, tada $38+4=42$.

\textbf{Atsakymas: 42.}

\textbf{2 pavyzdys.} Apskaičiuok: $53-20$.

Atimame 2 dešimtis iš 5 dešimčių: $50-20=30$. Vienetai lieka 3, todėl $53-20=33$.

\textbf{Atsakymas: 33.}

\textbf{3 pavyzdys.} Knygoje buvo 36 lipdukai. Vaikas priklijavo 9 lipdukus. Kiek lipdukų liko?

Reikia atimti: $36-9$. Patogu atimti 6 ir dar 3: $36-6=30$, $30-3=27$.

\textbf{Atsakymas: liko 27 lipdukai.}

\subsection{Ryšys su kitomis temomis}

Sudėtis ir atimtis remiasi skaičių iki 100 ir vietinės vertės supratimu iš temų \ref{sec:skaiciai-iki-100-ir-veiksmai} ir \ref{sec:skaiciai-iki-1000-ir-vietine-verte}. Šie veiksmai vėliau naudojami daugyboje ir dalyboje \ref{sec:daugybos-ir-dalybos-pradzia}, reiškiniuose \ref{sec:reiskiniai-ir-raidiniai-zymejimai} ir aritmetinių veiksmų savybėse \ref{sec:aritmetiniai-veiksmai-ir-savybes}.

\section{Daugybos ir dalybos pradžia}\label{sec:daugybos-ir-dalybos-pradzia}

\subsection{Apibrėžtys}

\begin{apibreztis}
Daugyba yra vienodų grupių sudėjimas. Pavyzdžiui, $3\cdot 4$ reiškia 3 grupes po 4.
\end{apibreztis}

\begin{apibreztis}
Dalyba yra dalijimas į lygias dalis arba ieškojimas, kiek lygių grupių galima sudaryti.
\end{apibreztis}

\begin{apibreztis}
Sandauga yra daugybos rezultatas, o dalmuo yra dalybos rezultatas.
\end{apibreztis}

\subsection{Teoremos ir savybės}

\begin{savybe}
Daugyba gali būti kartotinė sudėtis: $4\cdot 3=3+3+3+3$.
\end{savybe}

\begin{savybe}
Dauginant iš 1, skaičius nepasikeičia: $1\cdot 8=8$.
\end{savybe}

\begin{savybe}
Dalijant į lygias dalis, kiekvienoje dalyje turi būti vienodas daiktų skaičius.
\end{savybe}

\begin{pastaba}
Pirmoje ir antroje klasėje daugybą ir dalybą geriausia vaizduoti daiktų grupėmis, eilėmis, taškais arba piešiniais.
\end{pastaba}

\subsection{Taisyklės ir algoritmai}

\begin{enumerate}
\item Daugybos uždavinį paversk vienodų grupių sudėtimi.
\item Suskaičiuok, kiek grupių yra ir kiek daiktų yra kiekvienoje grupėje.
\item Dalybos uždavinyje paskirstyk daiktus po lygiai arba sudaryk vienodo dydžio grupes.
\item Atsakymą patikrink: daugyba ir dalyba yra susiję veiksmai.
\end{enumerate}

\subsection{Iliustruojantys pavyzdžiai}

\textbf{1 pavyzdys.} Ant 4 lėkščių yra po 3 obuolius. Kiek obuolių iš viso?

Tai yra 4 grupės po 3: $3+3+3+3=12$. Galime rašyti $4\cdot 3=12$.

\textbf{Atsakymas: 12 obuolių.}

\textbf{2 pavyzdys.} 10 saldainių padalyta po lygiai 2 vaikams. Kiek saldainių gavo kiekvienas?

Dalijame 10 į 2 lygias dalis. Kiekvienoje dalyje bus po 5, nes $5+5=10$.

\textbf{Atsakymas: kiekvienas gavo po 5 saldainius.}

\textbf{3 pavyzdys.} Kiek yra $5\cdot 2$?

$5\cdot 2$ reiškia 5 grupes po 2. Sudedame: $2+2+2+2+2=10$.

\textbf{Atsakymas: 10.}

\subsection{Ryšys su kitomis temomis}

Daugybos pradžia remiasi sudėtimi iš temos \ref{sec:sudetis-ir-atimtis}. Dalyba padeda suprasti paprastąsias trupmenas kaip dalis temoje \ref{sec:paprastosios-trupmenos-kaip-dalys}. Vėliau šie veiksmai gilinami temose \ref{sec:daugyba-ir-dalyba-stulpeliu} ir \ref{sec:aritmetiniai-veiksmai-ir-savybes}.

\section{Paprastosios trupmenos kaip dalys}\label{sec:paprastosios-trupmenos-kaip-dalys}

\subsection{Apibrėžtys}

\begin{apibreztis}
Trupmena parodo vieną arba kelias lygias visumos dalis.
\end{apibreztis}

\begin{apibreztis}
Pusė yra viena iš dviejų lygių dalių. Ji žymima $\frac{1}{2}$.
\end{apibreztis}

\begin{apibreztis}
Ketvirtadalis yra viena iš keturių lygių dalių. Ji žymima $\frac{1}{4}$.
\end{apibreztis}

\subsection{Teoremos ir savybės}

\begin{savybe}
Visuma gali būti dalijama į lygias dalis: pusės turi būti vienodos, ketvirtadaliai taip pat turi būti vienodi.
\end{savybe}

\begin{savybe}
Dvi pusės sudaro vieną visumą: $\frac{1}{2}+\frac{1}{2}=1$.
\end{savybe}

\begin{savybe}
Keturios ketvirtosios dalys sudaro vieną visumą: $\frac{1}{4}+\frac{1}{4}+\frac{1}{4}+\frac{1}{4}=1$.
\end{savybe}

\begin{pastaba}
Trupmeną lengviau suprasti, kai matome dalijamą pyragą, lapą, juostelę, laikrodžio ratą arba daiktų grupę.
\end{pastaba}

\subsection{Taisyklės ir algoritmai}

\begin{enumerate}
\item Įsitikink, kad visuma padalyta į lygias dalis.
\item Suskaičiuok, į kiek lygių dalių padalyta visuma.
\item Nuspalvink arba paimk tiek dalių, kiek prašoma.
\item Trupmenos pavadinimą siek su dalių skaičiumi: 2 dalys yra pusės, 4 dalys yra ketvirtadaliai.
\end{enumerate}

\subsection{Iliustruojantys pavyzdžiai}

\textbf{1 pavyzdys.} Obuolys perpjautas į 2 lygias dalis. Kaip vadinama viena dalis?

Kai visuma padalyta į 2 lygias dalis, viena dalis vadinama puse. Ji žymima $\frac{1}{2}$.

\textbf{Atsakymas: viena dalis yra pusė, arba $\frac{1}{2}$.}

\textbf{2 pavyzdys.} Šokoladas padalytas į 4 lygias dalis. Vaikas suvalgė 1 dalį. Kokią šokolado dalį jis suvalgė?

Visuma padalyta į 4 lygias dalis, paimta 1 dalis. Tai yra viena ketvirtoji: $\frac{1}{4}$.

\textbf{Atsakymas: vaikas suvalgė $\frac{1}{4}$ šokolado.}

\textbf{3 pavyzdys.} Lapas padalytas į 2 lygias dalis ir abi dalys nuspalvintos. Kokia lapo dalis nuspalvinta?

Nuspalvintos abi pusės. Dvi pusės sudaro visą lapą: $\frac{1}{2}+\frac{1}{2}=1$.

\textbf{Atsakymas: nuspalvintas visas lapas.}

\subsection{Ryšys su kitomis temomis}

Trupmenos kaip dalys siejasi su dalyba iš temos \ref{sec:daugybos-ir-dalybos-pradzia} ir matavimais temoje \ref{sec:matavimai-ilgis-mase-turis-laikas-pinigai}. Vėliau trupmenos plečiamos temose \ref{sec:trupmenos-ir-desimtaines-trupmenos-pradzia} ir \ref{sec:paprastosios-ir-desimtaines-trupmenos}.

\section{Dėsningumai ir sekos}\label{sec:desningumai-ir-sekos}

\subsection{Apibrėžtys}

\begin{apibreztis}
Dėsningumas yra pasikartojanti arba aiškiai kintanti tvarka.
\end{apibreztis}

\begin{apibreztis}
Seka yra daiktų, ženklų arba skaičių eilė, sudaryta pagal tam tikrą taisyklę.
\end{apibreztis}

\begin{apibreztis}
Sekos narys yra vienas sekos elementas.
\end{apibreztis}

\subsection{Teoremos ir savybės}

\begin{savybe}
Jeigu seka turi taisyklę, pagal ją galima rasti kitą sekos narį.
\end{savybe}

\begin{savybe}
Skaičių sekoje nariai gali didėti, mažėti arba kartotis.
\end{savybe}

\begin{savybe}
Daiktų sekoje gali kartotis spalvos, formos, dydžiai arba kryptys.
\end{savybe}

\subsection{Taisyklės ir algoritmai}

\begin{enumerate}
\item Pažiūrėk į pirmus kelis sekos narius.
\item Paklausk: kas kartojasi arba kas keičiasi?
\item Pasakyk taisyklę žodžiais.
\item Pagal taisyklę įrašyk trūkstamą arba kitą narį.
\end{enumerate}

\subsection{Iliustruojantys pavyzdžiai}

\textbf{1 pavyzdys.} Tęsk seką: 2, 4, 6, 8, \ldots

Kiekvienas skaičius padidėja 2. Po 8 pridedame 2: $8+2=10$.

\textbf{Atsakymas: kitas skaičius yra 10.}

\textbf{2 pavyzdys.} Rask trūkstamą narį: raudona, mėlyna, raudona, mėlyna, \ldots

Spalvos kartojasi poromis: raudona, mėlyna. Po mėlynos vėl eina raudona.

\textbf{Atsakymas: trūkstama spalva yra raudona.}

\textbf{3 pavyzdys.} Tęsk seką: 15, 14, 13, 12, \ldots

Skaičiai mažėja po 1. Po 12 eina $12-1=11$.

\textbf{Atsakymas: kitas skaičius yra 11.}

\subsection{Ryšys su kitomis temomis}

Dėsningumai padeda greičiau skaičiuoti temose \ref{sec:skaiciai-iki-100-ir-veiksmai} ir \ref{sec:sudetis-ir-atimtis}. Jie susiję su lygybėmis ir nežinomuoju \ref{sec:lygybes-nelygybes-ir-nezinomasis}, o vėliau plečiami temose \ref{sec:sekos-lenteles-ir-koordinates} ir \ref{sec:sekos-ir-matematine-indukcija}.

\section{Lygybės, nelygybės ir nežinomasis}\label{sec:lygybes-nelygybes-ir-nezinomasis}

\subsection{Apibrėžtys}

\begin{apibreztis}
Lygybė rodo, kad kairėje ir dešinėje pusėje yra tiek pat. Lygybės ženklas yra $=$.
\end{apibreztis}

\begin{apibreztis}
Nelygybė rodo, kad viena pusė yra didesnė arba mažesnė. Naudojami ženklai $>$ ir $<$.
\end{apibreztis}

\begin{apibreztis}
Nežinomasis yra trūkstamas skaičius, kurį reikia rasti. Pirmose klasėse jis gali būti žymimas langeliu, tašku arba raide.
\end{apibreztis}

\subsection{Teoremos ir savybės}

\begin{savybe}
Lygybėje abi pusės turi turėti vienodą vertę: $6+2=8$.
\end{savybe}

\begin{savybe}
Jeigu viena pusė turi daugiau daiktų, naudojamas ženklas $>$, o jei mažiau, naudojamas ženklas $<$.
\end{savybe}

\begin{savybe}
Nežinomąjį galima rasti bandant, skaičiuojant pirmyn arba naudojant priešingą veiksmą.
\end{savybe}

\begin{pastaba}
Ženklas $<$ atsiveria į didesnį skaičių taip, kaip plačiau atverta pusė žiūri į didesnę krūvelę.
\end{pastaba}

\subsection{Taisyklės ir algoritmai}

\begin{enumerate}
\item Apskaičiuok abi puses, jeigu jose yra veiksmų.
\item Palygink gautus skaičius.
\item Pasirink ženklą $=$, $>$ arba $<$.
\item Jeigu yra nežinomasis, klausk, kiek trūksta arba kiek reikia atimti.
\end{enumerate}

\subsection{Iliustruojantys pavyzdžiai}

\textbf{1 pavyzdys.} Įrašyk ženklą: $7+2\ \square\ 10$.

Pirmiausia apskaičiuojame $7+2=9$. Lyginame 9 ir 10. Kadangi 9 mažiau už 10, rašome $<$.

\textbf{Atsakymas: $7+2<10$.}

\textbf{2 pavyzdys.} Rask nežinomąjį: $\square+4=9$.

Klausiame, kiek trūksta nuo 4 iki 9. Skaičiuojame: 5, 6, 7, 8, 9. Tai 5 žingsniai, todėl nežinomasis yra 5.

\textbf{Atsakymas: $\square=5$.}

\textbf{3 pavyzdys.} Ar teisinga lygybė $12-3=8$?

Apskaičiuojame: $12-3=9$. Kadangi 9 nėra 8, lygybė neteisinga.

\textbf{Atsakymas: lygybė neteisinga.}

\subsection{Ryšys su kitomis temomis}

Lygybės ir nelygybės remiasi skaičių lyginimu iš temų \ref{sec:skaiciai-iki-100-ir-veiksmai} ir \ref{sec:skaiciai-iki-1000-ir-vietine-verte}. Nežinomasis paruošia lygčių mokymąsi temose \ref{sec:paprastos-lygtys-ir-nelygybes} ir \ref{sec:lygtys-ir-nelygybes-su-vienu-nezinomuoju}.

\section{Plokštumos figūros ir kūnai}\label{sec:plokstumos-figuros-ir-kunai}

\subsection{Apibrėžtys}

\begin{apibreztis}
Plokštumos figūra yra figūra, kurią galima nupiešti lape. Pavyzdžiai: trikampis, kvadratas, stačiakampis, skritulys.
\end{apibreztis}

\begin{apibreztis}
Kūnas yra erdvinis daiktas, kurį galima paimti į rankas arba įsivaizduoti erdvėje. Pavyzdžiai: kubas, rutulys, cilindras.
\end{apibreztis}

\begin{apibreztis}
Kraštinė yra figūros šonas, o viršūnė yra vieta, kur susitinka kraštinės.
\end{apibreztis}

\subsection{Teoremos ir savybės}

\begin{savybe}
Trikampis turi 3 kraštines ir 3 viršūnes.
\end{savybe}

\begin{savybe}
Kvadratas turi 4 lygias kraštines ir 4 viršūnes.
\end{savybe}

\begin{savybe}
Skritulys neturi kraštinių ir kampų.
\end{savybe}

\begin{pastaba}
Kūnus galima sieti su kasdieniais daiktais: kamuolys primena rutulį, dėžutė primena kubą arba stačiakampį gretasienį.
\end{pastaba}

\subsection{Taisyklės ir algoritmai}

\begin{enumerate}
\item Atpažindamas figūrą, suskaičiuok jos kraštines ir viršūnes.
\item Palygink figūros kraštinių ilgius.
\item Kūną atpažink pagal jo paviršių: ar jis rieda, ar turi plokščias sieneles.
\item Piešdamas figūrą, pradėk nuo viršūnių arba nuo aiškios kreivos linijos.
\end{enumerate}

\subsection{Iliustruojantys pavyzdžiai}

\textbf{1 pavyzdys.} Figūra turi 3 kraštines. Kokia tai figūra?

Figūra su 3 kraštinėmis yra trikampis.

\textbf{Atsakymas: trikampis.}

\textbf{2 pavyzdys.} Kuo skiriasi kvadratas ir stačiakampis?

Abu turi po 4 kraštines ir 4 viršūnes. Kvadrato visos kraštinės lygios, o stačiakampio priešingos kraštinės lygios.

\textbf{Atsakymas: kvadrato visos kraštinės lygios, o stačiakampio lygios priešingos kraštinės.}

\textbf{3 pavyzdys.} Kurį kūną primena kamuolys?

Kamuolys yra apvalus ir gali riedėti į visas puses. Tokį kūną vadiname rutuliu.

\textbf{Atsakymas: kamuolys primena rutulį.}

\subsection{Ryšys su kitomis temomis}

Figūros ir kūnai siejasi su matavimais temoje \ref{sec:matavimai-ilgis-mase-turis-laikas-pinigai}. Vėliau mokomasi kampų, daugiakampių ir erdvinių kūnų temose \ref{sec:kampai-tieses-ir-daugiakampiai}, \ref{sec:kampai-trikampiai-ir-keturkampiai} ir \ref{sec:erdviniai-kunai-ir-ju-pjuviai}.

\section{Matavimai: ilgis, masė, tūris, laikas, pinigai}\label{sec:matavimai-ilgis-mase-turis-laikas-pinigai}

\subsection{Apibrėžtys}

\begin{apibreztis}
Matavimas yra dydžio palyginimas su pasirinktu matavimo vienetu.
\end{apibreztis}

\begin{apibreztis}
Ilgis rodo, koks daiktas ilgas, trumpas, aukštas arba žemas. Dažni vienetai yra centimetras ir metras.
\end{apibreztis}

\begin{apibreztis}
Masė rodo, koks daiktas sunkus arba lengvas. Dažni vienetai yra gramas ir kilogramas.
\end{apibreztis}

\begin{apibreztis}
Tūris rodo, kiek telpa inde. Laikas rodo įvykių trukmę arba kada jie vyksta. Pinigai naudojami kainai ir mokėjimui.
\end{apibreztis}

\subsection{Teoremos ir savybės}

\begin{savybe}
Norint tiksliai matuoti, reikia pradėti nuo nulio žymos.
\end{savybe}

\begin{savybe}
1 metras yra 100 centimetrų, o 1 euras yra 100 centų.
\end{savybe}

\begin{savybe}
Laikrodyje valanda turi 60 minučių.
\end{savybe}

\begin{pastaba}
Pirmose klasėse svarbu mokėti pasirinkti tinkamą vienetą: sąsiuvinio ilgį matuojame centimetrais, kambario ilgį metrais.
\end{pastaba}

\subsection{Taisyklės ir algoritmai}

\begin{enumerate}
\item Pasirink, ką matuosi: ilgį, masę, tūrį, laiką arba pinigus.
\item Pasirink tinkamą matavimo priemonę: liniuotę, svarstykles, matavimo indą, laikrodį arba monetas.
\item Pradėk matuoti nuo teisingos vietos, dažnai nuo nulio.
\item Užrašyk skaičių ir matavimo vienetą.
\end{enumerate}

\subsection{Iliustruojantys pavyzdžiai}

\textbf{1 pavyzdys.} Pieštukas prasideda ties 0 cm ir baigiasi ties 12 cm. Koks pieštuko ilgis?

Ilgį randame nuo pradžios iki pabaigos. Nuo 0 cm iki 12 cm yra 12 cm.

\textbf{Atsakymas: pieštuko ilgis yra 12 cm.}

\textbf{2 pavyzdys.} Obuolys kainuoja 30 centų, o kriaušė 40 centų. Kiek kainuoja abu vaisiai?

Sudedame kainas: $30+40=70$. Abu vaisiai kainuoja 70 centų.

\textbf{Atsakymas: 70 centų.}

\textbf{3 pavyzdys.} Pamoka prasidėjo 9 valandą ir truko 1 valandą. Kada pamoka baigėsi?

Prie 9 valandos pridedame 1 valandą: $9+1=10$.

\textbf{Atsakymas: pamoka baigėsi 10 valandą.}

\subsection{Ryšys su kitomis temomis}

Matavimai naudoja skaičius ir veiksmus iš temų \ref{sec:skaiciai-iki-100-ir-veiksmai}, \ref{sec:skaiciai-iki-1000-ir-vietine-verte} ir \ref{sec:sudetis-ir-atimtis}. Figūrų matavimas vėliau veda į perimetrą, plotą ir tūrį temose \ref{sec:perimetras-plotas-turis-ir-masto-pradzia} ir \ref{sec:plotas-pavirsius-ir-turis}.

\section{Duomenų rinkimas ir diagramos}\label{sec:duomenu-rinkimas-ir-diagramos}

\subsection{Apibrėžtys}

\begin{apibreztis}
Duomenys yra surinkta informacija, pavyzdžiui, vaikų pasirinktos spalvos arba suskaičiuoti daiktai.
\end{apibreztis}

\begin{apibreztis}
Lentelė yra tvarkingas duomenų užrašymas eilutėmis ir stulpeliais.
\end{apibreztis}

\begin{apibreztis}
Diagrama yra duomenų vaizdas, padedantis greitai pamatyti, ko daugiau, ko mažiau ir kiek yra iš viso.
\end{apibreztis}

\subsection{Teoremos ir savybės}

\begin{savybe}
Tą patį duomenų rinkinį galima parodyti žodžiais, lentele arba diagrama.
\end{savybe}

\begin{savybe}
Diagramos stulpelio aukštis arba paveikslėlių skaičius rodo, kiek kartų pasirinktas daiktas ar atsakymas.
\end{savybe}

\begin{savybe}
Duomenis lengviau lyginti, kai jie surašyti tvarkingai.
\end{savybe}

\subsection{Taisyklės ir algoritmai}

\begin{enumerate}
\item Nuspręsk, kokį klausimą tirsi.
\item Surink atsakymus arba suskaičiuok daiktus.
\item Užrašyk duomenis lentelėje.
\item Nubraižyk diagramą ir pagal ją atsakyk į klausimus.
\end{enumerate}

\subsection{Iliustruojantys pavyzdžiai}

\textbf{1 pavyzdys.} Klasėje 6 vaikai pasirinko obuolį, 4 vaikai pasirinko kriaušę. Kurių vaisių pasirinkta daugiau?

Lyginame skaičius 6 ir 4. Kadangi $6>4$, obuolių pasirinkta daugiau.

\textbf{Atsakymas: daugiau pasirinkta obuolių.}

\textbf{2 pavyzdys.} Lentelėje parašyta: raudoni pieštukai -- 5, mėlyni pieštukai -- 7. Kiek pieštukų iš viso?

Sudedame duomenis: $5+7=12$.

\textbf{Atsakymas: iš viso 12 pieštukų.}

\textbf{3 pavyzdys.} Paveikslėlių diagramoje yra 3 saulės ir 8 debesys. Ko pavaizduota mažiau?

Lyginame 3 ir 8. Kadangi $3<8$, saulių pavaizduota mažiau.

\textbf{Atsakymas: mažiau pavaizduota saulių.}

\subsection{Ryšys su kitomis temomis}

Duomenims reikia skaičiavimo, lyginimo ir sudėties iš temų \ref{sec:skaiciai-iki-100-ir-veiksmai}, \ref{sec:sudetis-ir-atimtis} ir \ref{sec:lygybes-nelygybes-ir-nezinomasis}. Vėliau diagramos plečiamos temose \ref{sec:duomenu-lenteles-ir-diagramos} ir \ref{sec:statistiniai-duomenys-ir-vidurkiai}.

\section{Tikimybės žodžiai ir įvykiai}\label{sec:tikimybes-zodziai-ir-ivykiai}

\subsection{Apibrėžtys}

\begin{apibreztis}
Įvykis yra tai, kas gali atsitikti. Pavyzdžiui, rytoj lis arba iš maišelio ištrauksime raudoną kaladėlę.
\end{apibreztis}

\begin{apibreztis}
Tikimybės žodžiai padeda kalbėti apie įvykius: neįmanoma, galima, tikėtina, tikra.
\end{apibreztis}

\begin{apibreztis}
Tikras įvykis būtinai įvyks, o neįmanomas įvykis negali įvykti.
\end{apibreztis}

\subsection{Teoremos ir savybės}

\begin{savybe}
Jeigu maišelyje yra tik raudonos kaladėlės, ištraukti raudoną kaladėlę yra tikras įvykis.
\end{savybe}

\begin{savybe}
Jeigu maišelyje nėra mėlynų kaladėlių, ištraukti mėlyną kaladėlę yra neįmanomas įvykis.
\end{savybe}

\begin{savybe}
Jeigu įvykis gali atsitikti, bet nebūtinai atsitiks, jį vadiname galimu įvykiu.
\end{savybe}

\begin{pastaba}
Pirmose klasėse tikimybę pakanka aptarti žodžiais ir pavyzdžiais iš kasdienio gyvenimo.
\end{pastaba}

\subsection{Taisyklės ir algoritmai}

\begin{enumerate}
\item Perskaityk situaciją ir išsiaiškink, kas gali atsitikti.
\item Pažiūrėk, kokie daiktai arba pasirinkimai galimi.
\item Nuspręsk, ar įvykis tikras, galimas, tikėtinas ar neįmanomas.
\item Paaiškink pasirinkimą paprastais žodžiais.
\end{enumerate}

\subsection{Iliustruojantys pavyzdžiai}

\textbf{1 pavyzdys.} Maišelyje yra 5 geltonos kaladėlės ir nėra kitų spalvų. Ar galima ištraukti geltoną kaladėlę?

Visos kaladėlės yra geltonos, todėl ištraukti geltoną kaladėlę tikrai pavyks.

\textbf{Atsakymas: tai tikras įvykis.}

\textbf{2 pavyzdys.} Dėžėje yra tik raudoni ir žali kamuoliukai. Ar galima ištraukti mėlyną kamuoliuką?

Dėžėje mėlynų kamuoliukų nėra. Todėl mėlyno kamuoliuko ištraukti negalima.

\textbf{Atsakymas: tai neįmanomas įvykis.}

\textbf{3 pavyzdys.} Lauke yra debesų. Ar gali šiandien lyti?

Kai yra debesų, lietus gali prasidėti, bet nebūtinai prasidės. Todėl sakome, kad tai galimas įvykis.

\textbf{Atsakymas: tai galimas įvykis.}

\subsection{Ryšys su kitomis temomis}

Tikimybės žodžiai siejasi su duomenų rinkimu temoje \ref{sec:duomenu-rinkimas-ir-diagramos}, nes stebėjimai padeda kalbėti apie tai, kas dažniau ar rečiau pasitaiko. Vėliau ši tema tęsiama \ref{sec:paprastos-tikimybes}, \ref{sec:atsitiktiniai-bandymas-ir-tikimybe} ir \ref{sec:tikimybe-ir-kombinatorika}.


\section{Finansiniai skaičiavimai: eurai ir centai}\label{sec:finansiniai-skaiciavimai-eurai-ir-centai}

\subsection{Apibrėžtys}
\begin{apibreztis}
Tema apima pinigų atpažinimą, kainų palyginimą, sumų sudarymą banknotais ir monetomis. Ji nagrinėjama kaip atskira mokymo(si) turinio dalis pagal \emph{matematikaPrograma.pdf}.
\end{apibreztis}
\begin{apibreztis}
Matematinis modelis šioje temoje yra skaičių, simbolių, brėžinių, lentelių arba diagramų sistema, kuria aprašoma reali arba teorinė situacija.
\end{apibreztis}

\subsection{Teoremos ir savybės}
\begin{savybe}
Jeigu uždavinio sąlyga tiksliai perkeliama į matematinį modelį, tai modelio rezultatas gali būti interpretuojamas pradinėje situacijoje.
\end{savybe}
\begin{pastaba}
Pagrindimas remiasi nagrinėjamų objektų apibrėžtimis ir tuo, kad kiekvienas veiksmas išlaiko pradinės situacijos prasmę. Taikant taisyklę svarbu remtis konkrečiais modeliais, brėžiniais arba lentele.
\end{pastaba}

\subsection{Taisyklės ir algoritmai}
\begin{enumerate}[label=\arabic*.]
\item Išskirkite duotus dydžius, sąvokas arba objektus.
\item Pasirinkite tinkamą užrašą: skaitinį reiškinį, formulę, lentelę, brėžinį, diagramą arba lygtį.
\item Atlikite veiksmus nuosekliai, kiekviename žingsnyje išlaikydami vienetus ir sąlygas.
\item Patikrinkite, ar gautas rezultatas atitinka uždavinio prasmę.
\end{enumerate}

\subsection{Iliustruojantys pavyzdžiai}
\begin{enumerate}[label=\arabic*.]
\item Pritaikykime temą paprastam skaičiavimui: pradinė reikšmė yra \(12\), pokytis yra \(3\).
\begin{align}
12+3&=15.
\end{align}
\textbf{Atsakymas: \(15\).}
\item Palyginkime du galimus rezultatus \(18\) ir \(24\).
\begin{align}
24-18&=6,\\
24&>18.
\end{align}
\textbf{Atsakymas: antras rezultatas didesnis \(6\) vienetais.}
\item Sudarykime ir išspręskime paprastą modelį, kai nežinomas dydis \(x\) su \(5\) sudaro \(20\).
\begin{align}
x+5&=20,\\
x&=15.
\end{align}
\textbf{Atsakymas: \(x=15\).}
\end{enumerate}

\subsection{Ryšys su kitomis temomis}
Ši tema papildo to paties koncentro skaičių, modelių, geometrijos arba duomenų temas ir padeda nuosekliai pereiti prie aukštesnių klasių uždavinių, kuriuose reikia derinti kelias matematikos sritis.


\section{Algoritmai ir programavimas: komandos}\label{sec:algoritmai-ir-programavimas-komandos}

\subsection{Apibrėžtys}
\begin{apibreztis}
Tema apima komandų sekas, žodžius „ne“, „arba“, „ir“ ir paprasčiausią programos vykdymą. Ji nagrinėjama kaip atskira mokymo(si) turinio dalis pagal \emph{matematikaPrograma.pdf}.
\end{apibreztis}
\begin{apibreztis}
Matematinis modelis šioje temoje yra skaičių, simbolių, brėžinių, lentelių arba diagramų sistema, kuria aprašoma reali arba teorinė situacija.
\end{apibreztis}

\subsection{Teoremos ir savybės}
\begin{savybe}
Jeigu uždavinio sąlyga tiksliai perkeliama į matematinį modelį, tai modelio rezultatas gali būti interpretuojamas pradinėje situacijoje.
\end{savybe}
\begin{pastaba}
Pagrindimas remiasi nagrinėjamų objektų apibrėžtimis ir tuo, kad kiekvienas veiksmas išlaiko pradinės situacijos prasmę. Taikant taisyklę svarbu remtis konkrečiais modeliais, brėžiniais arba lentele.
\end{pastaba}

\subsection{Taisyklės ir algoritmai}
\begin{enumerate}[label=\arabic*.]
\item Išskirkite duotus dydžius, sąvokas arba objektus.
\item Pasirinkite tinkamą užrašą: skaitinį reiškinį, formulę, lentelę, brėžinį, diagramą arba lygtį.
\item Atlikite veiksmus nuosekliai, kiekviename žingsnyje išlaikydami vienetus ir sąlygas.
\item Patikrinkite, ar gautas rezultatas atitinka uždavinio prasmę.
\end{enumerate}

\subsection{Iliustruojantys pavyzdžiai}
\begin{enumerate}[label=\arabic*.]
\item Pritaikykime temą paprastam skaičiavimui: pradinė reikšmė yra \(12\), pokytis yra \(3\).
\begin{align}
12+3&=15.
\end{align}
\textbf{Atsakymas: \(15\).}
\item Palyginkime du galimus rezultatus \(18\) ir \(24\).
\begin{align}
24-18&=6,\\
24&>18.
\end{align}
\textbf{Atsakymas: antras rezultatas didesnis \(6\) vienetais.}
\item Sudarykime ir išspręskime paprastą modelį, kai nežinomas dydis \(x\) su \(5\) sudaro \(20\).
\begin{align}
x+5&=20,\\
x&=15.
\end{align}
\textbf{Atsakymas: \(x=15\).}
\end{enumerate}

\subsection{Ryšys su kitomis temomis}
Ši tema papildo to paties koncentro skaičių, modelių, geometrijos arba duomenų temas ir padeda nuosekliai pereiti prie aukštesnių klasių uždavinių, kuriuose reikia derinti kelias matematikos sritis.


\section{Matavimo skalės ir vienetai}\label{sec:matavimo-skales-ir-vienetai-pradzia}

\subsection{Apibrėžtys}
\begin{apibreztis}
Tema apima masės, laiko, ilgio, atstumo, temperatūros, talpos ir ploto matavimą praktinėse situacijose. Ji nagrinėjama kaip atskira mokymo(si) turinio dalis pagal \emph{matematikaPrograma.pdf}.
\end{apibreztis}
\begin{apibreztis}
Matematinis modelis šioje temoje yra skaičių, simbolių, brėžinių, lentelių arba diagramų sistema, kuria aprašoma reali arba teorinė situacija.
\end{apibreztis}

\subsection{Teoremos ir savybės}
\begin{savybe}
Jeigu uždavinio sąlyga tiksliai perkeliama į matematinį modelį, tai modelio rezultatas gali būti interpretuojamas pradinėje situacijoje.
\end{savybe}
\begin{pastaba}
Pagrindimas remiasi nagrinėjamų objektų apibrėžtimis ir tuo, kad kiekvienas veiksmas išlaiko pradinės situacijos prasmę. Taikant taisyklę svarbu remtis konkrečiais modeliais, brėžiniais arba lentele.
\end{pastaba}

\subsection{Taisyklės ir algoritmai}
\begin{enumerate}[label=\arabic*.]
\item Išskirkite duotus dydžius, sąvokas arba objektus.
\item Pasirinkite tinkamą užrašą: skaitinį reiškinį, formulę, lentelę, brėžinį, diagramą arba lygtį.
\item Atlikite veiksmus nuosekliai, kiekviename žingsnyje išlaikydami vienetus ir sąlygas.
\item Patikrinkite, ar gautas rezultatas atitinka uždavinio prasmę.
\end{enumerate}

\subsection{Iliustruojantys pavyzdžiai}
\begin{enumerate}[label=\arabic*.]
\item Pritaikykime temą paprastam skaičiavimui: pradinė reikšmė yra \(12\), pokytis yra \(3\).
\begin{align}
12+3&=15.
\end{align}
\textbf{Atsakymas: \(15\).}
\item Palyginkime du galimus rezultatus \(18\) ir \(24\).
\begin{align}
24-18&=6,\\
24&>18.
\end{align}
\textbf{Atsakymas: antras rezultatas didesnis \(6\) vienetais.}
\item Sudarykime ir išspręskime paprastą modelį, kai nežinomas dydis \(x\) su \(5\) sudaro \(20\).
\begin{align}
x+5&=20,\\
x&=15.
\end{align}
\textbf{Atsakymas: \(x=15\).}
\end{enumerate}

\subsection{Ryšys su kitomis temomis}
Ši tema papildo to paties koncentro skaičių, modelių, geometrijos arba duomenų temas ir padeda nuosekliai pereiti prie aukštesnių klasių uždavinių, kuriuose reikia derinti kelias matematikos sritis.


\section{Konstravimas ir transformacijos: judėjimas plane}\label{sec:konstravimas-ir-transformacijos-judejimas-plane}

\subsection{Apibrėžtys}
\begin{apibreztis}
Tema apima objektų padėtį, judėjimo kryptis, posūkius, simetriją ir planus languotame popieriuje. Ji nagrinėjama kaip atskira mokymo(si) turinio dalis pagal \emph{matematikaPrograma.pdf}.
\end{apibreztis}
\begin{apibreztis}
Matematinis modelis šioje temoje yra skaičių, simbolių, brėžinių, lentelių arba diagramų sistema, kuria aprašoma reali arba teorinė situacija.
\end{apibreztis}

\subsection{Teoremos ir savybės}
\begin{savybe}
Jeigu uždavinio sąlyga tiksliai perkeliama į matematinį modelį, tai modelio rezultatas gali būti interpretuojamas pradinėje situacijoje.
\end{savybe}
\begin{pastaba}
Pagrindimas remiasi nagrinėjamų objektų apibrėžtimis ir tuo, kad kiekvienas veiksmas išlaiko pradinės situacijos prasmę. Taikant taisyklę svarbu remtis konkrečiais modeliais, brėžiniais arba lentele.
\end{pastaba}

\subsection{Taisyklės ir algoritmai}
\begin{enumerate}[label=\arabic*.]
\item Išskirkite duotus dydžius, sąvokas arba objektus.
\item Pasirinkite tinkamą užrašą: skaitinį reiškinį, formulę, lentelę, brėžinį, diagramą arba lygtį.
\item Atlikite veiksmus nuosekliai, kiekviename žingsnyje išlaikydami vienetus ir sąlygas.
\item Patikrinkite, ar gautas rezultatas atitinka uždavinio prasmę.
\end{enumerate}

\subsection{Iliustruojantys pavyzdžiai}
\begin{enumerate}[label=\arabic*.]
\item Pritaikykime temą paprastam skaičiavimui: pradinė reikšmė yra \(12\), pokytis yra \(3\).
\begin{align}
12+3&=15.
\end{align}
\textbf{Atsakymas: \(15\).}
\item Palyginkime du galimus rezultatus \(18\) ir \(24\).
\begin{align}
24-18&=6,\\
24&>18.
\end{align}
\textbf{Atsakymas: antras rezultatas didesnis \(6\) vienetais.}
\item Sudarykime ir išspręskime paprastą modelį, kai nežinomas dydis \(x\) su \(5\) sudaro \(20\).
\begin{align}
x+5&=20,\\
x&=15.
\end{align}
\textbf{Atsakymas: \(x=15\).}
\end{enumerate}

\subsection{Ryšys su kitomis temomis}
Ši tema papildo to paties koncentro skaičių, modelių, geometrijos arba duomenų temas ir padeda nuosekliai pereiti prie aukštesnių klasių uždavinių, kuriuose reikia derinti kelias matematikos sritis.
